% figs.tex —— 由 ch4/make_figs.py 產生，請勿手動修改。
% 重新產生：python ch4/make_figs.py（第 4 章，動畫 W401–W408）

\newcommand{\FourFigMirror}{%
\begin{tikzpicture}[x=1.1cm,y=0.48cm,>=Stealth,dot/.style={circle,inner sep=1.2pt}]
\begin{scope}[xshift=0.0cm]
\draw[gray!25] (0,-2) grid (4,3);
\draw[cWrong,very thick] (0,0) -- (4.25,0) node[right,font=\scriptsize] {懸崖 $0$};
\foreach \y in {-2,-1,0,1,2,3} \node[left,font=\scriptsize] at (0,\y) {$\y$};
\foreach \k in {0,...,4} \node[below,font=\scriptsize] at (\k,-2) {$\k$};
\node[font=\scriptsize] at (2,-3.3) {步數 $k$};
\node[font=\scriptsize] at (2,3.7) {$1\to0\to1\to2\to3\ \leftrightarrow\ -1\to0\to1\to2\to3$};
\draw[cTail,very thick,dashed] (0,-1) -- (1,0);
\draw[cWalk,very thick] (0,1) -- (1,0);
\draw[cWalk,very thick] (1,0) -- (2,1) -- (3,2) -- (4,3);
\draw[cTail,thick,densely dashed] (1,0) -- (2,1) -- (3,2) -- (4,3);
\node[dot,fill=cWalk] at (0,1) {};
\node[dot,fill=cWalk] at (1,0) {};
\node[dot,fill=cWalk] at (2,1) {};
\node[dot,fill=cWalk] at (3,2) {};
\node[dot,fill=cWalk] at (4,3) {};
\node[dot,fill=cTail] at (0,-1) {};
\node[above right,font=\scriptsize,cWalk] at (0.06,1) {起點 $x_0=1$};
\node[below right,font=\scriptsize,cTail] at (0.06,-1) {鏡像點 $-x_0=-1$};
\node[circle,draw=cWrong,thick,inner sep=2.6pt] at (1,0) {};
\node[below right,font=\scriptsize,cWrong] at (1,0) {第一次碰牆};
\end{scope}
\begin{scope}[xshift=7.4cm]
\draw[gray!25] (0,-2) grid (4,3);
\draw[cWrong,very thick] (0,0) -- (4.25,0) node[right,font=\scriptsize] {懸崖 $0$};
\foreach \y in {-2,-1,0,1,2,3} \node[left,font=\scriptsize] at (0,\y) {$\y$};
\foreach \k in {0,...,4} \node[below,font=\scriptsize] at (\k,-2) {$\k$};
\node[font=\scriptsize] at (2,-3.3) {步數 $k$};
\node[font=\scriptsize] at (2,3.7) {$1\to2\to1\to0\to1\ \leftrightarrow\ -1\to-2\to-1\to0\to1$};
\draw[cTail,very thick,dashed] (0,-1) -- (1,-2) -- (2,-1) -- (3,0);
\draw[cWalk,very thick] (0,1) -- (1,2) -- (2,1) -- (3,0);
\draw[cWalk,very thick] (3,0) -- (4,1);
\draw[cTail,thick,densely dashed] (3,0) -- (4,1);
\node[dot,fill=cWalk] at (0,1) {};
\node[dot,fill=cWalk] at (1,2) {};
\node[dot,fill=cWalk] at (2,1) {};
\node[dot,fill=cWalk] at (3,0) {};
\node[dot,fill=cWalk] at (4,1) {};
\node[dot,fill=cTail] at (0,-1) {};
\node[dot,fill=cTail] at (1,-2) {};
\node[dot,fill=cTail] at (2,-1) {};
\node[anchor=north west,inner sep=1pt,font=\scriptsize,cWalk] at (0.06,0.9) {起點 $x_0=1$};
\node[anchor=south west,inner sep=1pt,font=\scriptsize,cTail] at (0.06,-0.9) {鏡像點 $-x_0=-1$};
\node[circle,draw=cWrong,thick,inner sep=2.6pt] at (3,0) {};
\node[below right,font=\scriptsize,cWrong] at (3,0) {第一次碰牆};
\end{scope}
\end{tikzpicture}}

\newcommand{\FourTabSurvive}{%
\begin{tabular}{@{}lccc@{}}
\toprule
起點 $x_0=3$、牆在 $0$ & $N=10$ & $N=100$ & $N=1000$\\
\midrule
一維：精確存活機率 & $21/32=0.65625$ & 0.2356 & 0.07557\\
一維：$x_0\sqrt{2/\pi N}$ & 0.7569 & 0.2394 & 0.07569\\
sc 格子、吸收平面：精確 & -- & 0.3958 & 0.1305\\
三維：$\sqrt{6x_0^2/\pi N}$（(4.5) 後） & -- & 0.4146 & 0.1311\\
\midrule
一維 (4.8)：牆吃掉的（$\div2^N$）精確 & $\frac{1067}{512}=2.084$ & 7.499 & \\
一維 (4.8)：$\sqrt{2N/\pi}$ & 2.523 & 7.979 & \\
\bottomrule
\end{tabular}}

\newcommand{\FourFigDecay}{%
\begin{tikzpicture}[x=0.06cm,y=0.54cm,>=Stealth]
\draw[->] (0,-6) -- (214,-6) node[right,font=\scriptsize] {$N$};
\draw[->] (0,-6) -- (0,0.7) node[above,font=\scriptsize] {存活比例（對數刻度）};
\foreach \n in {0,50,100,150,200} {\draw (\n,-6.06) -- (\n,-5.94); \node[below,font=\scriptsize] at (\n,-6) {$\n$};}
\foreach \e in {0,-1,-2,-3,-4,-5,-6} {\draw (-1.5,\e) -- (1.5,\e); \node[left,font=\scriptsize] at (-1.5,\e) {$10^{\e}$};}
\draw[gray!25] (0,-1) -- (200,-1) (0,-2) -- (200,-2) (0,-3) -- (200,-3) (0,-4) -- (200,-4) (0,-5) -- (200,-5);
\draw[cBar,thick] (0,-0.0526) -- (200.00,-4.4114);
\draw[cExact,very thick,dashed] (0,-0.0912) -- (200.00,-4.3775);
\draw[cWrong,thick,dash dot] (0,-0.0912) -- (137.85,-6.0000);
\fill[cBar] (0,0.0000) circle (1.7pt);
\fill[cBar] (10,-0.2699) circle (1.7pt);
\fill[cBar] (20,-0.4882) circle (1.7pt);
\fill[cBar] (30,-0.7062) circle (1.7pt);
\fill[cBar] (40,-0.9241) circle (1.7pt);
\fill[cBar] (50,-1.1420) circle (1.7pt);
\fill[cBar] (60,-1.3600) circle (1.7pt);
\fill[cBar] (70,-1.5779) circle (1.7pt);
\fill[cBar] (80,-1.7959) circle (1.7pt);
\fill[cBar] (90,-2.0138) circle (1.7pt);
\fill[cBar] (100,-2.2317) circle (1.7pt);
\fill[cBar] (110,-2.4497) circle (1.7pt);
\fill[cBar] (120,-2.6676) circle (1.7pt);
\fill[cBar] (130,-2.8855) circle (1.7pt);
\fill[cBar] (140,-3.1035) circle (1.7pt);
\fill[cBar] (150,-3.3214) circle (1.7pt);
\fill[cBar] (160,-3.5394) circle (1.7pt);
\fill[cBar] (170,-3.7573) circle (1.7pt);
\fill[cBar] (180,-3.9752) circle (1.7pt);
\fill[cBar] (190,-4.1932) circle (1.7pt);
\fill[cBar] (200,-4.4111) circle (1.7pt);
\node[right,font=\scriptsize,cBar] at (88,-0.15) {點：精確計數 $\div 2^N(L-1)$；藍線：格子最低模態 $\propto\cos^N(\pi/L)$};
\node[right,font=\scriptsize,cExact] at (88,-0.8) {綠虛線：更正 $\frac{8}{\pi^2}e^{-\pi^2N/2L^2}$（和藍線幾乎重合）};
\node[right,font=\scriptsize,cWrong] at (88,-1.45) {紅點劃線：書上 (4.24) $\frac{8}{\pi^2}e^{-\pi^2N/L^2}$（兩倍陡）};
\end{tikzpicture}}

\newcommand{\FourTabRates}{%
\begin{tabular}{@{}lccc@{}}
\toprule
牆距 $L$ & 格子 $-\ln\cos(\pi/L)$ & 更正 $\pi^2/2L^2$ & {\color{cWrong}書上 $\pi^2/L^2$}\\
\midrule
10 & 0.05018 & 0.04935 & 0.0987\\
20 & 0.01239 & 0.01234 & 0.0247\\
\midrule 牆距加倍，衰減變慢 & 4.05 倍 & 4 倍 & 4 倍\\
\bottomrule
\end{tabular}}

\newcommand{\FourTabShapes}{%
{\small\setlength{\tabcolsep}{4pt}\begin{tabular}{@{}lcccc@{}}
\toprule
區域（吸收牆） & 短時間總數 & 長時間總數 & 前因子 $\div$ 體積 & 每步衰減率\\
\midrule
線段 $0<x<L$（一維） & $2^NL$ & $2^N\frac{8L}{\pi^2}e^{-\pi^2N/2L^2}$ & $8/\pi^2=0.811$ & $\pi^2/2L^2$\\
立方體（邊長 $L$，練習 4.5） & $6^NL^3$ & $6^N\frac{512L^3}{\pi^6}e^{-\pi^2N/2L^2}$ & $512/\pi^6=0.533$ & $\pi^2/2L^2$\\
球（半徑 $R$） & $6^N\cdot\frac43\pi R^3$ & $6^N\frac{8R^3}{\pi}e^{-\pi^2N/6R^2}$ & $6/\pi^2=0.608$ & $\pi^2/6R^2$\\
\bottomrule
\end{tabular}}}

\newcommand{\FourFigTent}{%
\begin{tikzpicture}[x=1.75cm,y=0.8cm,>=Stealth]
\draw[->] (-0.3,0) -- (5.6,0) node[right,font=\scriptsize] {$x$};
\foreach \x in {0,...,5} \node[below,font=\scriptsize] at (\x,0) {$\x$};
\fill[cBar!20,draw=cBar] (0.8,0) rectangle (1.2,1.2000);
\node[font=\scriptsize,cBar] at (1,0.6000) {$\frac{6}{5}$};
\fill[cBar!20,draw=cBar] (1.8,0) rectangle (2.2,2.4000);
\node[font=\scriptsize,cBar] at (2,1.2000) {$\frac{12}{5}$};
\fill[cBar!20,draw=cBar] (2.8,0) rectangle (3.2,1.6000);
\node[font=\scriptsize,cBar] at (3,0.8000) {$\frac{8}{5}$};
\fill[cBar!20,draw=cBar] (3.8,0) rectangle (4.2,0.8000);
\node[font=\scriptsize,cBar] at (4,0.4000) {$\frac{4}{5}$};
\draw[cExact,thick] (0,0) -- (2,2.4000) -- (5,0);
\draw[cWrong,line width=2pt] (0,0) -- (0,2.9) (5,0) -- (5,2.9);
\node[above,font=\scriptsize,cWrong] at (0,2.9) {吸收牆};
\node[above,font=\scriptsize,cWrong] at (5,2.9) {吸收牆};
\draw[->,cWalk,very thick] (2,3.45) -- (2,2.5200);
\node[above,font=\scriptsize,cWalk] at (2,3.45) {每步在 $x_0=2$ 放 1 人（$s_0=1$）};
\node[font=\scriptsize,cExact] at (0.5,1.45) {斜率 $\frac65$};
\node[font=\scriptsize,cExact] at (3.75,1.6) {斜率 $-\frac45$};
\node[right,font=\scriptsize,cExact] at (2.12,2.7500) {斜率跳 $-s_0/\beta=-2$};
\draw[->,cWrong,thick] (0.85,-0.75) -- (0.05,-0.75);
\node[below,font=\scriptsize,cWrong] at (0.45,-0.75) {流進左牆 $w(1)/2=\frac{3}{5}$};
\draw[->,cWrong,thick] (4.15,-0.75) -- (4.95,-0.75);
\node[below,font=\scriptsize,cWrong] at (4.55,-0.75) {流進右牆 $w(4)/2=\frac{2}{5}$};
\end{tikzpicture}}

\newcommand{\FourTabGreen}{%
{\small\begin{tabular}{@{}lcccccccc@{}}
\toprule
$n$（沿 $x$ 軸） & 0 & 1 & 2 & 3 & 4 & 5 & 10 & $\infty$\\
\midrule
$w=G(z_c;n\hat x)$ & 1.5164 & 0.5164 & 0.2573 & 0.1653 & 0.1217 & 0.0966 & 0.04787 & $\to0$\\
$n\,G$ & -- & 0.5164 & 0.5147 & 0.4958 & 0.4869 & 0.4830 & 0.4787 & $\frac{3}{2\pi}=0.4775$\\
\bottomrule
\end{tabular}}}

\newcommand{\FourTabBP}{%
\begin{tabular}{@{}lccccc@{}}
\toprule
受器數 $n$（$r_0=1\,\mu$m、$a=1$ nm） & 100 & 1000 & 3142 & $10^4$ & $10^5$\\
\midrule
$J/J_{\max}=\frac{na}{na+\pi r_0}$ & 0.031 & 0.241 & 0.500 & 0.761 & 0.970\\
受器占表面 $\frac{na^2}{4r_0^2}$ & 0.0025\% & 0.025\% & 0.079\% & 0.25\% & 2.5\%\\
\bottomrule
\end{tabular}}

\newcommand{\FourTabErrata}{%
\begin{tabular}{@{}lll@{}}
\toprule
位置 & 書上 & 應為\\
\midrule
p.~75，(4.12)(4.13) & {\color{cWrong}$\sqrt{2/L}\,\sin qx_1$} & {\color{cExact}$\sqrt{2/L}\,\sin qx_2$}\\
p.~75，(4.15)〔注〕 & {\color{cWrong}$\sum_{n=0}^{\infty}$} & {\color{cExact}$\sum_{n=1}^{\infty}$（$n=0$ 項為 0）}\\
p.~75 末行 & {\color{cWrong}$0\le z\le L$} & {\color{cExact}$0\le x\le L$}\\
p.~77，(4.22) & {\color{cWrong}$\frac14\sum_{k=0}^\infty\frac1{k^2}$、$\frac34\sum_{n=0}^\infty\frac1{n^2}$} & {\color{cExact}和從 1 起（$k=0$ 項發散）}\\
p.~77，(4.24) & {\color{cWrong}$e^{-\pi^2\Delta^2N/L^2}$} & {\color{cExact}$e^{-\pi^2\Delta^2N/2L^2}$}\\
p.~80，(4.31) & {\color{cWrong}$\frac{8\pi}{R}\cdots\frac2R\frac{j_0j_0}{1-2z+z\Delta^2k^2}$} & {\color{cExact}$\frac{8\pi}{R}\cdots k_{n_0}^2\frac{j_0j_0}{1-6z+z\Delta^2k^2}$}\\
p.~81，(4.35) & {\color{cWrong}$e^{-\Delta^2\pi^2N/R^2}$} & {\color{cExact}$e^{-\Delta^2\pi^2n^2N/6R^2}$}\\
p.~81，(4.37) & {\color{cWrong}$(6e^{-\Delta^2\pi^2/R^2})^N$} & {\color{cExact}$(6e^{-\Delta^2\pi^2/6R^2})^N$}\\
p.~83，(4.40) & {\color{cWrong}$\frac{a^2}{3}\nabla^2w$} & {\color{cExact}$\frac{a^2}{6}\nabla^2w$}\\
p.~83，(4.41) & {\color{cWrong}$\frac{3}{a^2}S$} & {\color{cExact}$\frac{6}{a^2}S$}\\
p.~87，(4.58) 前 & {\color{cWrong}wall in the $x,z$ plane} & {\color{cExact}$y,z$ 平面（$x=0$）}\\
p.~91，(4.64) 前 & {\color{cWrong}$\rho=-4\pi\nabla^2\phi$} & {\color{cExact}$\nabla^2\phi=-4\pi\rho$}\\
p.~92 & {\color{cWrong}$na\ll r_0$} & {\color{cExact}$na^2\ll r_0^2$（$a\ll r_0/\sqrt n$）}\\
p.~93 & {\color{cWrong}$a\ll r/\sqrt n$} & {\color{cExact}$a\ll r_0/\sqrt n$}\\
p.~93，(4:S-1) 後 & {\color{cWrong}starting at $y$ and ending at $y$} & {\color{cExact}starting at $x$ and ending at $y$}\\
p.~93，(4:S-2) & {\color{cWrong}$G(z;0,y)=zG(z;x,\Delta)+\delta_{x,\Delta}$} & {\color{cExact}$G(z;x,0)=zG(z;x,\Delta)+\delta_{x,0}$}\\
\bottomrule
\end{tabular}}


\section{邊界條件、穩定態與靜電類比（第 4 章）}

\begin{keybox}[先看這裡]
前三章的醉漢走在沒有盡頭的街上。這一章在街上挖了懸崖：\textbf{吸收牆}＝「踩到就消失」的格點。問題變成「還有多少人活著」，
以及「源源不絕地放人進來，等很久以後人群怎麼分布」。三件工具：\textbf{吸收牆 $\to$ 鏡像}（碰過牆的路線＝從鏡像點出發的路線，相減就好）；
\textbf{有限區域 $\to$ 模態}（正弦波或 $j_0=\sin x/x$ 取代平面波，長時間活著的比例 $\propto e^{-\pi^2N/2L^2}$）；
\textbf{穩定態 $\to$ 靜電學}（源＋漫步＝Poisson 方程，吸收牆＝接地導體，被吸收的流量＝電容）。總整理在 4.5 節。
\textbf{記號}：步長 $\Delta=a=1$；$N$＝步數；$\mu$＝每步可走的方向數（簡單立方 sc 為 6）；一面牆在 $x=0$，起點 $x_0\ge1$，鏡像點 $-x_0$；
兩面牆在 $0$ 與 $L$（內點 $1,\dots,L-1$）；穩定態每步在源放 $s_0$ 個走路人，格子上的濃度記 $w$、連續的記 $c$；擴散係數 $\beta$（＝$D$）。

\textbf{全章的主例}：第 1 章的醉漢，「0 是懸崖」。從 $x_0=1$ 走 4 步，16 條路線只有 6 條沒掉下去（4.1.1 節）；
懸崖在 0 和 6、從 3 出發，走 10 步的 1024 條路線裡活著 324 條（4.1.2 節）；走廊 $0..5$、每步在 2 放進一個人，
穩定後各格的平均人數是「帳篷」$\frac65,\frac{12}5,\frac85,\frac45$（4.2.1 節）。
\begin{center}
\small
\begin{tabular}{@{}lll@{\qquad}lll@{}}
\toprule
節 & 內容 & 動畫 & 節 & 內容 & 動畫\\
\midrule
4.1.1 & 鏡像走路人 & W401 & 4.2.1 & 穩定態、Fick 定律、帳篷 & W405\\
4.1.2 & 兩面吸收牆：正弦模態 & W402 & 4.2.2--4.2.3 & 點源的 $1/r$、牆旁的點源 & W406\\
4.1.3 & 短時間與長時間 & W403 & 4.2.4--4.2.5 & 吸收球、電容 & W407\\
4.1.4 & 球與立方體 & W404 & 4.2.6、4.3 & 受器、吸收邊界、回顧 & W408\\
\bottomrule
\end{tabular}
\end{center}
數字一律由 \texttt{codes/utils/walk\_bc.py} 重算，圖表由 \texttt{ch4/make\_figs.py} 產生並先做數字檢查；逐式校對見 \texttt{ref/md/ch04\_formulas.md}。
連續播放：\texttt{Ch4Demo}（W401--W408）。
\end{keybox}

% =====================================================================
\subsection{吸收牆與有限區域}
\bk{pp.~69--82}\quad 本節都在「數活著的路線」：一面牆用鏡像，兩面牆用正弦模態，再搬到球與立方體裡。

\subsubsection{鏡像走路人}
\bk{pp.~69--73}\ \anim{W401}

\textbf{直觀}：醉漢從 $x_0\ge1$ 出發，0 是懸崖。要數「$N$ 步都沒踩到 0」的路線，笨方法是一條一條檢查；聰明的方法是先數全部，
再扣掉碰過牆的——而碰過牆的路線可以一一對應到「從鏡像點 $-x_0$ 出發」的路線。

\textbf{最小的例子}：$x_0=1$、$N=4$，共 $2^4=16$ 條路線。碰過 0 的有 10 條；活著的 6 條，終點在 1、3、5 的各有 2、3、1 條。
16 條裡終點 $>0$ 的有 11 條，其中 5 條碰過牆。這 5 條正好對應「從 $-1$ 出發、終點 $>0$」的 5 條：把\textbf{第一次碰牆之前}的那段對 0 鏡射，
之後兩條路線重合（圖 \ref{four-fig-mirror}）；反過來，從 $-1$ 走到正的地方一定要經過 0，所以也能鏡射回來。於是（$C(N;d)$＝沒有牆、位移 $d$ 的走法數）
\[C'(N;x,x_0)=C(N;x-x_0)-C(N;x+x_0)\qquad(x=1{:}\ 6-4=2,\ \ x=3{:}\ 4-1=3,\ \ x=5{:}\ 1-0=1),\]
共 $6=11-5$ 條。這就是書上 (4.1) 的一維版。
\begin{figure}[H]
\centering
\FourFigMirror
\caption{書上 Fig.~4.1 的樣子（$x_0=1$、$N=4$）。橘色是碰過牆的路線，藍色虛線是把第一次碰牆之前的那段鏡射後、從 $-1$ 出發的路線；
碰牆之後兩條重合。這樣的配對在 16 條裡共有 5 對}\label{four-fig-mirror}
\end{figure}

\textbf{書上的式子}：三維、牆是平面 $x=0$，起點 $\vec r_0=(x_0,y,z)$、鏡像點 $\vec r_I=(-x_0,y,z)$：
\begin{gather*}
C'(N;\vec r,\vec r_0)=C(N;\vec r-\vec r_0)-C(N;\vec r-\vec r_I),\tag{4.1}\\
C'(N;\vec r,\vec r_0)=\mu^N\Big(\frac{2\pi N}{3}\Big)^{-3/2}\Big[e^{-3|\vec r-\vec r_0|^2/2N}-e^{-3|\vec r-\vec r_I|^2/2N}\Big].\tag{4.2}
\end{gather*}
對 $y,z$ 積分（終點落在平面 $x=x_F$ 上）得 $\mu^N(\frac{2\pi N}{3})^{-1/2}\big[e^{-3(x_F-x_0)^2/2N}-e^{-3(x_F+x_0)^2/2N}\big]$（(4.3)）。
再對 $x_F\in(0,\infty)$ 積分：第一個高斯變成 $\int_{-x_0}^\infty$、第二個變成 $\int_{x_0}^\infty$，相減只剩中間一段：
\[\Gamma(N;x_0)=\mu^N\Big(\frac{2\pi N}{3}\Big)^{-1/2}\int_{-x_0}^{x_0}e^{-3x^2/2N}dx\ \approx\ \mu^N\Big(\frac{2\pi N}{3}\Big)^{-1/2}\cdot2x_0\quad(x_0^2\ll N).\tag{4.4c, 4.5}\]
所以活著的比例是 $\sqrt{6x_0^2/\pi N}$：隨 $N^{-1/2}$ 慢慢趨近 0。（〔注〕p.~72 說「$x_0^2\gg2N$ 時牆沒有影響」：指數是 $3x_0^2/2N$，
條件寫成 $x_0^2\gg2N/3$ 較準，反正只是數量級。）

\textbf{一維的對照}：把 $C'$ 對終點 $x>0$ 加總，$\sum_{x>0}C(N;x-x_0)$ 數的是「位移 $S_N>-x_0$」、$\sum_{x>0}C(N;x+x_0)$ 數的是「$S_N>x_0$」，
所以存活機率＝$P(1-x_0\le S_N\le x_0)$（$S_N$ 是沒有牆的位移）。例：$x_0=3$、$N=10$，$S_{10}\in\{-2,0,2\}$，存活機率 $(210+252+210)/1024=21/32$。
$N$ 很大時它 $\approx x_0\sqrt{2/\pi N}$（$2x_0\times$ 高斯密度；奇偶性的因子 2 剛好被「只有一半的整數和 $N$ 同奇偶」抵掉）。
三維 sc 的答案大 $\sqrt3$ 倍：$x$ 座標只在 $1/3$ 的步數裡移動，三維走 $N$ 步碰平面，等於一維走 $N/3$ 步碰點。
\begin{table}[H]
\centering
\small
\FourTabSurvive
\caption{起點 $x_0=3$ 的存活機率：精確值與漸近式（一維與 sc 都是 $\propto N^{-1/2}$，一維終究一定掉下去）。下兩列是一維版的 (4.8)：
所有起點加起來被牆吃掉多少（$N=10$ 時漸近式還不準）}\label{four-tab-survive}
\end{table}

\textbf{牆吃掉多少}（(4.6)--(4.8)）：把起點 $x_0$ 也加起來，用 $\int_0^\infty dx_0\int_{x_0}^\infty e^{-3x^2/2N}dx=\int_0^\infty xe^{-3x^2/2N}dx=\frac N3$（交換積分次序）：
\[\text{牆吃掉的走法數}=2\mu^N\Big(\frac{2\pi N}{3}\Big)^{-1/2}\frac N3=\mu^N\Big(\frac{2N}{3\pi}\Big)^{1/2}.\tag{4.8}\]
讀法：離牆 $\sqrt N$ 之內出發的人大多活不久，更遠的幾乎不受影響。一維版是 $\sum_{x_0\ge1}(1-\text{存活機率})=E[\max_{k\le N}S_k]\approx\sqrt{2N/\pi}$
（從 $x_0$ 出發會掉下去 $\Leftrightarrow$ 某一步 $-S_k\ge x_0$）：$N=10$ 精確 $1067/512=2.084$、漸近 2.523；$N=100$ 精確 7.499、漸近 7.979。

\subsubsection{兩面吸收牆：正弦模態}
\bk{pp.~73--76}\ \anim{W402}

\textbf{例子}：懸崖在 0 和 6，醉漢從正中間 3 出發。在每個格點上寫「走到這裡、還沒掉下去」的走法數；每一步把左右鄰居相加，牆上永遠是 0
（$x{:}n$＝在 $x$ 有 $n$ 條）：
\[N=0{:}\ \{3{:}1\},\ \ N=1{:}\ \{2{:}1,\,4{:}1\},\ \ N=2{:}\ \{1{:}1,\,3{:}2,\,5{:}1\},\ \ N=3{:}\ \{2{:}3,\,4{:}3\},\ \ N=4{:}\ \{1{:}3,\,3{:}6,\,5{:}3\}.\]
從 $N=2$ 到 $N=4$ 每格都乘 3，之後也是每兩步乘 3：$N=10$ 時 $x=1,\dots,5$ 依序是 $81,\ 0,\ 162,\ 0,\ 81$，共 324 條（全部 1024 條，活著約 0.316）。

\textbf{為什麼是 3？}「把左右鄰居相加」是一個 $5\times5$ 的轉移矩陣。它的特徵向量是 $\sin(m\pi x/6)$（$m=1,\dots,5$）——在 $x=0$ 與 $x=6$ 自動是 0，
正好符合牆的條件——因為 $\sin q(x+1)+\sin q(x-1)=2\cos q\,\sin qx$（$q=m\pi/L$），
特徵值 $2\cos(m\pi/6)=\sqrt3,\ 1,\ 0,\ -1,\ -\sqrt3$。把起點拆成模態，每個模態每走一步乘上自己的特徵值：
\[C_{\rm box}(N;x_0,x)=\sum_{m=1}^{L-1}\frac2L\sin\frac{m\pi x_0}{L}\sin\frac{m\pi x}{L}\Big(2\cos\frac{m\pi}{L}\Big)^N.\]
從中間 $x_0=3$ 出發，$m=2,4$ 的 $\sin(m\pi/2)=0$ 沒有份；$m=3$ 的特徵值是 0；只剩 $m=1,5$，而 $(\pm\sqrt3)^{10}=243$：
$x=3$ 是 $\frac26(243+243)=162$；$x=1$ 用 $\sin\frac\pi6=\sin\frac{5\pi}6=\frac12$，得 $\frac26(\frac12\cdot243+\frac12\cdot243)=81$；$x=2$ 的兩項正負抵消，是 0。

\textbf{第三種算法：無窮多個鏡像}（練習 4.3）。一面牆用一個鏡像點；兩面牆要「鏡像的鏡像……」：正的鏡像在 $2kL+x_0$（$\dots,-9,3,15,\dots$），
負的在 $2kL-x_0$（$\dots,-15,-3,9,\dots$）。$x=3$：$C(10;0)-C(10;6)-C(10;-6)=252-45-45=162$，其餘鏡像太遠、走不到。

\textbf{書上的式子}：連續版的模態是 $\psi(x,y,z)=e^{ik_yy+ik_zz}\sin qx$，$q=n\pi/L$、$n=1,2,3,\dots$（(4.9)(4.10)）。只看 $x$ 方向，生成函數滿足 (4.11)；把 $G$ 展開成 $\sqrt{2/L}\,\sin qx_1$ 的組合、係數記 $g(z;q)$，代進去比較係數得 (4.12)(4.13)：
\begin{gather*}
G(z;x_1,x_2)=z\big(G(z;x_1+\Delta,x_2)+G(z;x_1-\Delta,x_2)\big)+\delta_{x_1,x_2},\tag{4.11}\\
g(z;q)=2z\cos q\Delta\,g(z;q)+\sqrt{\frac2L}\sin qx_2,\qquad g(z;q)=\sqrt{\frac2L}\,\frac{\sin qx_2}{1-2z\cos q\Delta},\tag{4.12, 4.13}\\
G(z;x_1,x_2)=\frac2L\sum_{n=1}^{\infty}\frac{\sin\frac{n\pi x_1}{L}\sin\frac{n\pi x_2}{L}}{1-2z+z\Delta^2(\frac{n\pi}{L})^2}.\tag{4.15}
\end{gather*}
中間用了 (4.14)：$1-2z\cos q\Delta\approx(1-2z)+z\Delta^2q^2$（即 $2\cos q\Delta\approx2-\Delta^2q^2$）。所以 (4.15) 是上面格子版的連續極限；格子版取 $z^N$ 係數就是 $C_{\rm box}$。
\begin{errbox}[p.~75，(4.12)(4.13)(4.15) 與末行]
(4.12)(4.13) 的 $\sqrt{2/L}\,\sin qx_1$ 應為 $\sqrt{2/L}\,\sin qx_2$：遞迴移動的是 $x_1$，展開係數裡留下的是另一端 $x_2$（(4.14) 已經寫對）。
〔注〕(4.15) 的和從 $n=0$ 起：那一項本來就是 0，寫 $n=1$ 才和 (4.10) 一致。末行的 $0\le z\le L$ 應為 $0\le x\le L$。
\end{errbox}

\subsubsection{短時間像沒牆、長時間指數衰減}
\bk{pp.~76--79}\ \anim{W403}

\textbf{問}：兩面牆的距離加倍，「還活著的比例」衰減變慢幾倍？2 倍？4 倍？

\textbf{總數}：起點、終點都在 $(0,L)$ 裡加總。$\int_0^L\sin\frac{n\pi x}{L}dx=\frac{L}{n\pi}(1-\cos n\pi)$ 在偶數 $n$ 是 0、奇數 $n$ 是 $\frac{2L}{n\pi}$（(4.16)），
所以只剩奇數模態：$\frac{8L}{\pi^2}\sum_{m=0}^\infty\frac{1}{(2m+1)^2}\frac{1}{1-2z+z\Delta^2((2m+1)\pi/L)^2}$（(4.17)）。
取 $z^N$ 係數，$(2-\Delta^2q^2)^N\to2^Ne^{-N\Delta^2q^2/2}$（(4.18)）：
\[2^N\frac{8L}{\pi^2}\sum_{m=0}^\infty\frac{1}{(2m+1)^2}\exp\Big[-\frac{N\Delta^2}{2}\Big(\frac{(2m+1)\pi}{L}\Big)^2\Big].\tag{4.19}\]

\textbf{短時間}（$N\ll L^2$）：指數都約為 1，需要奇數平方倒數的和。由 $\sum_{n\ge1}\frac1{n^2}=\frac{\pi^2}{6}=1.6449$（(4.21)），偶數項是全部的 $\frac14$：
\[\sum_{m=0}^\infty\frac1{(2m+1)^2}=\sum_{n=1}^\infty\frac1{n^2}-\frac14\sum_{k=1}^\infty\frac1{k^2}=\frac34\cdot\frac{\pi^2}6=\frac{\pi^2}8=1.2337.\tag{4.22}\]
部分和 $1,\ 1.1111,\ 1.1511,\ 1.1715,\dots\to1.2337$。代回 (4.19)：$2^N\frac{8L}{\pi^2}\cdot\frac{\pi^2}{8}=2^NL$（(4.23)），和沒有牆一樣——每個起點各 $2^N$ 條
（格子上是 $2^N(L-1)$：$L-1$ 個內點）。
\begin{errbox}[p.~77，(4.22)]
第二、三行的 $\frac14\sum_{k=0}^\infty\frac1{k^2}$、$\frac34\sum_{n=0}^\infty\frac1{n^2}$：$k=0$ 那一項是 $1/0$，和應從 1 起。
\end{errbox}

\textbf{長時間}：$m\ge1$ 的項衰減得快得多（指數大 9 倍、25 倍……），只剩 $m=0$：$2^N\frac{8L}{\pi^2}e^{-\pi^2\Delta^2N/2L^2}$（(4.24)）。
\begin{errbox}[p.~77，(4.24)]
書上印成 $e^{-\pi^2\Delta^2N/L^2}$：(4.19) 的 $m=0$ 項有 $\frac12$，(4.24) 把它掉了。格子上的精確衰減率是每步 $-\ln\cos(\pi/L)$（表 \ref{four-tab-rates}）。
\end{errbox}
\begin{table}[H]
\centering
\small
\FourTabRates
\caption{每步的衰減率。牆距加倍，衰減慢 4 倍（格子 4.05 倍）；書上的指數每一格都大了 2 倍}\label{four-tab-rates}
\end{table}
格子上活著的比例由最低模態 $\frac{2}{L(L-1)}\cot^2\frac{\pi}{2L}\cos^N\frac\pi L$ 主導（$\sum_{x=1}^{L-1}\sin\frac{m\pi x}{L}=\cot\frac{m\pi}{2L}$，$m$ 奇，是 (4.16) 的格子版），
$L\to\infty$ 時前因子 $\to8/\pi^2=0.8106$。圖 \ref{four-fig-decay} 是 $L=10$ 的精確計數與三條直線。
\begin{figure}[H]
\centering
\FourFigDecay
\caption{$L=10$：活著的比例（所有起點平均）對 $N$ 的對數圖。長時間是一條直線；精確值在 $N=50$、100、200 是 0.072103、0.005865、$3.881\times10^{-5}$。
$N=100$ 時三條線是 0.005861（格子最低模態）、0.005830（更正）、$4.19\times10^{-5}$（書上）}\label{four-fig-decay}
\end{figure}

\textbf{揭曉}：衰減率 $\pi^2/2L^2\propto1/L^2$，牆距加倍、衰減慢 4 倍。每多走 $2L^2/\pi^2$ 步（$L=10$：20.26 步；格子上 19.93 步）活著的比例少 $e$ 倍。

\textbf{小例子再看一次}（$L=6$，所有起點加總）：$N=0,\dots,10$ 活著的路線有 5、8、14、24、42、72、126、216、378、648、1134 條。
$N=1$ 就少了 2 條（$5\times2=10$：兩端的人往外踏一步就掉下去）；之後每兩步剛好乘 3：奇數模態 $m=1,3,5$ 裡 $m=3$ 的特徵值是 0，剩下 $\pm\sqrt3$。
活著的比例 $\propto(\sqrt3/2)^N=\cos^N(\pi/6)$，每步衰減 $-\ln\cos(\pi/6)=0.1438$。

\textbf{熵與漏失}（(4.25)--(4.27)，Fig.~4.2）：書上的直觀說法。每走一步，每條路線分裂成 2 條；牆上的濃度是 0，濃度從中間的約 $n_w/L$ 降到牆邊，
斜率約 $\kappa n_w/L^2$（$\kappa$ 是形狀因子），流量 $J=-\alpha\,dn/dx$ 把這麼多漏進牆裡：
$n_w(N)=n_w(N-1)\times2\times(1-\frac{\alpha\kappa}{L^2})$，所以 $n_w(N)\approx2^Ne^{-\alpha\kappa N/L^2}$。量綱論證給不出常數；精確計算給 $\alpha\kappa=\pi^2/2=4.935$。

\subsubsection{關在球裡與立方體裡}
\bk{pp.~79--82}\ \anim{W404}

\textbf{問}：同樣的體積，關在球裡和關在立方體裡，哪一個活得久？

\textbf{球的模態}：在半徑 $R$ 的吸收球裡，模態是球 Bessel 函數乘球諧函數 $j_l(k_{n_l}r)Y_l^m(\theta,\phi)$，牆上 $j_l(k_{n_l}R)=0$，
分母是三維的 $1-6z+z\Delta^2k_{n_l}^2$（(4.28)(4.29)）。起點、終點都在球裡加總時，角度積分只留 $l=m=0$（(4.30)）；
$j_0(x)=\frac{\sin x}{x}$ 的零點是 $k_{n_0}=\frac{n_0\pi}{R}$（$k_{n_0}R=\pi,\ 2\pi,\ 3\pi,\dots$），徑向積分算得出來：
\begin{gather*}
\frac{8\pi}{R}\sum_{n_0}k_{n_0}^2\int_0^Rr_1^2dr_1\int_0^Rr_2^2dr_2\,\frac{j_0(k_{n_0}r_1)j_0(k_{n_0}r_2)}{1-6z+z\Delta^2k_{n_0}^2}\tag{4.31}\\
=\frac{8R^3}{\pi}\sum_n\frac1{n^2}\,\frac1{1-6z+z\Delta^2\pi^2n^2/R^2}.\tag{4.32--4.34}
\end{gather*}
權重 $\frac{8R^3}{\pi n^2}$ 的由來：$4\pi\int_0^Rr^2j_0(kr)dr=\pm\frac{4\pi R}{k^2}$，平方後除以正規化 $\int|j_0(kr)|^2d^3r=\frac{2\pi R}{k^2}$，得 $\frac{8\pi R}{k^2}=\frac{8R^3}{\pi n^2}$。
$R=1$ 時 $n=1$ 是 2.546、$n=2$ 是 0.637（程式另用數值積分核對）。
\begin{errbox}[p.~80，(4.31)]
書上印成 $\frac{8\pi}{R}\cdots\frac{2}{R}\frac{j_0j_0}{1-2z+z\Delta^2k^2}$：正規化給的是 $k^2$ 不是 $\frac2R$；三維的分母是 $1-6z$。照書上算 $R=1$、$n=1$ 的權重得 0.516，不是 2.546。
\end{errbox}

取 $z^N$ 係數，$(6-\Delta^2k^2)^N\approx6^Ne^{-N\Delta^2k^2/6}$：
\begin{gather*}
6^N\frac{8R^3}{\pi}\sum_{n=1}^\infty\frac1{n^2}e^{-\Delta^2\pi^2n^2N/6R^2}\ \xrightarrow{\ \text{短時間}\ }\ 6^N\frac{8R^3}{\pi}\cdot\frac{\pi^2}6=6^N\Big(\frac43\pi R^3\Big),\\
\xrightarrow{\ \text{長時間}\ }\ \frac{8R^3}{\pi}\Big(6e^{-\Delta^2\pi^2/6R^2}\Big)^N.\tag{4.35--4.37}
\end{gather*}
短時間：$R=1$ 時 $\frac43\pi=4.189$，正是球的體積——短時間牆不重要。長時間的指數裡的 $\frac16$ 來自三維：每一步只有 $\frac13$ 的機會沿某個軸走，
$\frac{k^2}{6}$ 是一維 $\frac{q^2}{2}$ 的三維版。
\begin{errbox}[p.~81，(4.35)(4.37)]
(4.35) 印成 $e^{-\Delta^2\pi^2N/R^2}$（少了 $n^2$ 與 $\frac16$）；(4.37) 印成 $(6e^{-\Delta^2\pi^2/R^2})^N$（少了 $\frac16$）。
格子核對：半徑 8 的格子球（2103 個格點）每步衰減 0.0242；更正 $\pi^2/(6\cdot64)=0.0257$；書上 $\pi^2/64=0.154$，差 6 倍。
\end{errbox}

\textbf{立方體}（練習 4.5）：邊長 $L$ 的立方體，模態是三個方向的正弦相乘，特徵值相加（$6-\Delta^2|\vec k|^2$），三個方向各取奇數模態：
$6^N(\frac{8L}{\pi^2})^3\sum_{n_i\,\text{奇}}(n_1n_2n_3)^{-2}e^{-\pi^2\Delta^2N(n_1^2+n_2^2+n_3^2)/6L^2}$。短時間 $(\frac{8L}{\pi^2})^3(\frac{\pi^2}8)^3=L^3$，得 $6^NL^3$；長時間只剩 $n_1=n_2=n_3=1$：$6^N\frac{512L^3}{\pi^6}e^{-\pi^2\Delta^2N/2L^2}$，衰減率和一維線段相同。
格子 $L=6$（$5^3=125$ 個內點）：$N=0$ 有 125 條，$N=1$ 有 $125\times6-6\times25=600$ 條（每個面上的 25 個格點各有一個方向會撞牆），$N=2$ 有 2970 條，逐步遞迴與模態和一致。
\begin{table}[H]
\centering
\FourTabShapes
\caption{三種吸收區域的比較。前因子 $\div$ 體積：線段 0.811、立方體 0.533、球 0.608。同體積時立方體與球的衰減率之比是 $3(3/4\pi)^{2/3}=1.155$}\label{four-tab-shapes}
\end{table}
\textbf{揭曉}：同體積 $L^3=\frac43\pi R^3$ 時，立方體的衰減率是球的 $3R^2/L^2=1.155$ 倍：\textbf{球最耐關}
（書外：同體積的區域裡球的最低模態最慢，Faber--Krahn 不等式）。形狀只改衰減率與前因子。

% =====================================================================
\subsection{穩定態與靜電類比}
\bk{pp.~82--93}\quad 改成\textbf{每一步都放人}（像一直噴香水）。等很久以後，放進來的和被牆吃掉的打平，濃度不再改變：穩定態。
主旋律：穩定態的濃度滿足 Poisson 方程，和靜電學一模一樣。

\subsubsection{源＋漫步＝Poisson 方程；Fick 定律}
\bk{pp.~82--86}\ \anim{W405}

\textbf{問}：走廊 $0..5$，兩端是懸崖；每一步都在 $x_0=2$ 放進一個醉漢。等很久以後，每一格平均站著幾個人？

\textbf{格子上的穩定態}：$W(N;\vec r)$＝第 $N$ 步時站在 $\vec r$ 的人數。每個人每步往 6 個鄰居之一走，源每步再放 $S(\vec r)$ 個：
\[W(N+1;\vec r)=\frac16\sum_{\vec\rho_i}W(N;\vec\rho_i)+S(\vec r),\qquad w(\vec r)=\frac16\sum_{\vec\rho_i}w(\vec\rho_i)+S(\vec r)\tag{4.38, 4.39}\]
（$\vec\rho_i$ 是 $\vec r$ 的 6 個鄰居；右式是穩定後 $W\to w$）。泰勒展開：每一對 $\pm\hat e$ 給 $2w+a^2\partial^2w$，所以
\[w(\vec r)=w(\vec r)+\frac{a^2}{6}\nabla^2w(\vec r)+S(\vec r),\qquad -\nabla^2w(\vec r)=\frac{6}{a^2}S(\vec r).\tag{4.40, 4.41}\]
這就是 Poisson 方程。數值核對：$[\frac16\sum_if(\vec r+a\hat e_i)-f(\vec r)]/(a^2\nabla^2f)\to0.1667$。
\begin{errbox}[p.~83，(4.40)(4.41)]
書上印成 $\frac{a^2}{3}\nabla^2w$、$-\nabla^2w=\frac{3}{a^2}S$：$\frac16\sum_{\pm x,\pm y,\pm z}w=w+\frac{a^2}{6}\nabla^2w$，應為 $\frac{a^2}{6}$、$\frac{6}{a^2}$。
同頁的「square lattice」應為簡單立方格子。
\end{errbox}

\textbf{連續版：Fick 定律}（(4.42)--(4.45)）：粒子守恆加上「流量往濃度低的地方」：
\[\frac{\partial c}{\partial t}=-\nabla\cdot\vec j+S,\qquad \vec j=-\beta\nabla c,\qquad \frac{\partial c}{\partial t}=\beta\nabla^2c+S,\qquad -\beta\nabla^2c=S.\tag{4.42--4.45}\]
和 (4.41) 比較，格子上 $\beta=a^2/(2d\,\tau)$：一維 $\frac12$、三維 $\frac16$（每步）。\textbf{對照靜電學}：濃度 $c$ $\leftrightarrow$ 電位 $\phi$；
源 $S$ $\leftrightarrow$ 電荷密度 $\rho$；流量 $\vec j=-\beta\nabla c$ $\leftrightarrow$ 電場 $\vec E=-\nabla\phi$；吸收牆 $c=0$ $\leftrightarrow$ 接地導體；
反射牆 $\hat n\cdot\nabla c=0$（書上註 4）$\leftrightarrow$ 沒有電場穿過的面；被吸收的總流量 $\leftrightarrow$ 導體上的電荷（電容）。

\textbf{例題：帳篷}（(4.46)--(4.49)）。一維、兩面吸收牆在 $0$ 與 $L$、源在 $x_0$：$\beta c''=-s_0\delta(x-x_0)$（(4.47)）。源以外 $c''=0$，所以是兩段直線，
在牆上是 0，在 $x_0$ 接起來、斜率跳 $c'(x_0^+)-c'(x_0^-)=-s_0/\beta$（(4.48)）：
\[c(x,x_0)=\begin{cases}\dfrac{s_0}{\beta}\dfrac{L-x_0}{L}\,x & x<x_0\\[2mm] \dfrac{s_0}{\beta}\dfrac{L-x}{L}\,x_0 & x>x_0\end{cases}\tag{4.49}\]
\textbf{算到底}：$L=5$、$x_0=2$、$s_0=1$、$\beta=\frac12$。(4.49) 給 $w=0,\ \frac65,\ \frac{12}5,\ \frac85,\ \frac45,\ 0$（1.2、2.4、1.6、0.8），
而且在格子上剛好精確（直線本來就滿足「等於左右鄰居的平均」）。在源上核對 (4.39) 的一維版：$\frac12(\frac65+\frac85)+1=\frac{12}5$。
斜率 $\frac65$、$-\frac45$，跳了 $-2=-s_0/\beta$。
\begin{figure}[H]
\centering
\FourFigTent
\caption{走廊 $0..5$、每步在 $x_0=2$ 放 1 人的穩定態（長條：格子精確解；綠線：(4.49)）。每步流進左牆 $\frac35$、右牆 $\frac25$，合起來正好是放進來的 1 人}\label{four-fig-tent}
\end{figure}

\textbf{練習 4.6（賭徒破產）}：穩定時，站在 1 的人每步有一半往左掉，所以每步流進左牆 $w(1)/2=\frac35$、右牆 $w(4)/2=\frac25$，和是 1。
換成一個人的說法：從 2 出發、先掉進左邊懸崖的機率是 $\frac35=\frac{L-x_0}{L}$（連續版：流進左牆 $\beta c'(0)=s_0\frac{L-x_0}{L}$）。這就是賭徒破產問題。
$L\to\infty$ 時機率 $\to1$（例如 $L=1000$：$\frac{998}{1000}$）：一維走不到無窮遠，和 4.1.1 節一致。

\subsubsection{\texorpdfstring{點源的 $1/r$}{點源的 1/r}}
\bk{pp.~86--87}\ \anim{W406}

\textbf{問}：在原點不斷放出走路人，穩定以後，離源 $r$ 處的濃度怎麼隨 $r$ 變？像高斯那樣急速下降嗎？

點源 $S=s_0\delta(\vec r)$ 的 Poisson 方程 $-\beta\nabla^2c=s_0\delta(\vec r)$ 就是點電荷的 Coulomb 問題，解是 $c(\vec r)=\frac{s_0}{4\pi\beta|\vec r|}$（(4.50)--(4.52)）：
只是 $1/r$，比高斯慢得多。\textbf{另一個看法}：穩定態＝過去每一刻放出的那一群人（各自是一個高斯）的總和，也就是把高斯對時間積分。令 $t=Tr^2$ 把 $r$ 提出來：
\[c_S(\vec r)\propto\int_0^\infty t^{-3/2}e^{-r^2/at}dt=\frac1r\int_0^\infty T^{-3/2}e^{-1/aT}dT=\frac Kr,\qquad K=\sqrt{\pi a}\tag{4.55, 4.56}\]
（再令 $u=1/aT$，$K=\sqrt a\,\Gamma(\frac12)$）。正規化的版本：$\int_0^\infty(4\pi Dt)^{-3/2}e^{-r^2/4Dt}dt=\frac1{4\pi Dr}$；$D=1$、$r=2$ 時是 $0.03979=\frac{1}{8\pi}$。

\textbf{格子上}（sc，每步在原點放 1 人）：穩定時站在 $\vec r$ 的人數＝「$N$ 步前放的人現在在 $\vec r$ 的機率」對 $N$ 求和＝第 3 章的 Green 函數在 $z_c=\frac16$：$w(\vec r)=G(z_c;\vec r)=\int_0^\infty e^{-t}I_{n_1}(\tfrac t3)I_{n_2}(\tfrac t3)I_{n_3}(\tfrac t3)dt$。原點 $w(0)=G(z_c;0)=1.516386$ 正是第 3 章的 Watson 積分。更正後的 (4.41) 是 $-\nabla^2w=6\delta$，所以遠處 $w\approx\frac{6}{4\pi r}=\frac{3}{2\pi r}$：
\begin{table}[H]
\centering
\FourTabGreen
\caption{sc 格子、每步在原點放 1 人的穩定濃度 $w$（沿 $x$ 軸）。$n\,G$ 趨近 $3/2\pi=0.4775$；照書上的 $3/a^2$ 會得 $3/4\pi=0.2387$，差 2 倍}\label{four-tab-green}
\end{table}
一個好記的核對：原點的 (4.39) 是 $w(0)=\frac16\cdot6\,w(1)+1$，所以 $G(1)=G(0)-1$：$1.5164-1=0.5164$。

\subsubsection{吸收牆旁的點源：偶極}
\bk{pp.~87--89}\ \anim{W406}

源在 $\vec r_s=\hat x\,l$、吸收牆是平面 $x=0$。和 4.1.1 節一樣放一個鏡像：在 $\vec r_i=-\hat x\,l$ 放強度相反的源（像接地導體平面前的點電荷）：
\begin{gather*}
c(\vec r)\propto\frac1{|\vec r-\vec r_s|}-\frac1{|\vec r-\vec r_i|},\tag{4.57}\\
c(x,y,z)\propto\frac1{\sqrt{(x-l)^2+y^2+z^2}}-\frac1{\sqrt{(x+l)^2+y^2+z^2}}\ \to\ \frac{2lx}{(x^2+y^2+z^2)^{3/2}}.\tag{4.58}
\end{gather*}
遠處是\textbf{偶極}，比 $1/r$ 衰減得快。
\begin{errbox}[p.~87，(4.58) 前]
書上說牆在「$x,z$ 平面」：源在 $x$ 軸上（$\vec r_s=\hat xl$、$\vec r_i=-\hat xl$），牆要垂直 $x$ 軸，應是 $y,z$ 平面（$x=0$）。
\end{errbox}
\textbf{有多少人逃到遠處？}穿過遠處平面 $x=\text{常數}$ 的總流量（$R$ 是平面上離 $x$ 軸的距離）：
\begin{gather*}
j_x\propto-\frac{\partial}{\partial x}\frac{2lx}{r^3}=2l\Big(\frac{3x^2}{r^5}-\frac1{r^3}\Big),\tag{4.59}\\
\int_0^\infty\Big[\frac{3x^2}{(x^2+R^2)^{5/2}}-\frac1{(x^2+R^2)^{3/2}}\Big]RdR=\Big[-\frac{x^2}{(x^2+w)^{3/2}}+\frac1{(x^2+w)^{1/2}}\Big]_0^\infty=0\tag{4.60}
\end{gather*}
（$w=R^2$）。不用遠場也能精確算：一個點源穿過一個無限大平面的流量正好是它的一半（平面張開半個立體角）。平面在源的外側（$X>l$）：源 $+\frac12$、鏡像 $-\frac12$，合計 0；
平面在牆與源之間（$0<X<l$）：$-1$——全部流向牆。\textbf{最後全被牆吃掉}：三維走路人的 $x$ 座標本身是一維漫步，終究會碰到牆（4.1.1 節）。

\subsubsection{吸收球：細胞長不大}
\bk{pp.~89--90}\ \anim{W407}

\textbf{問}：細胞靠擴散吃營養。細胞半徑加倍，吃進來的營養變幾倍？需要的呢？

半徑 $r_0$ 的吸收球，遠處濃度 $c_\infty$（書上記 $c_0$）。球外 $\nabla^2c=0$ 的球對稱解是 $A+B/r$，配上 $c(r_0)=0$、$c(\infty)=c_\infty$：
\[c(\vec r)=c_\infty\Big(1-\frac{r_0}{r}\Big),\qquad \vec j=-\hat r\frac{\beta c_\infty r_0}{r^2},\qquad \beta c_\infty\frac{r_0}{r_0^2}\times4\pi r_0^2=4\pi\beta c_\infty r_0.\tag{4.61--4.63}\]
流量 $4\pi\beta c_\infty r_0$ 穿過任何半徑的球面都一樣（$r_0=\beta=c_\infty=1$ 時是 $4\pi=12.566$），而且和 $r_0$ 成正比，不是和表面積 $r_0^2$。

\textbf{細胞的上限}：吃進 $\propto r_0$、需要 $\frac43\pi r_0^3q$（$q$＝每單位體積的消耗率）。半徑加倍，吃進 $\times2$、需要 $\times8$。令兩者相等得
$r_{\max}=\sqrt{3Dc_\infty/q}$。例（\textbf{假設值，只看數量級}，不是書上的數字）：$D=2\times10^{-9}\,\mathrm{m^2/s}$（氧在水中）、$c_\infty=0.25\,\mathrm{mol/m^3}$、
$q=1\,\mathrm{mol/(m^3\,s)}$，得 $r_{\max}\approx3.9\times10^{-5}$ m（約 39\,$\mu$m）。只靠擴散的細胞長不大。

\textbf{球外的點源}（Fig.~4.6）：源離球心 $R$。在 $\rho=r_0^2/R$ 放強度 $-r_0/R$ 的鏡像源，球面上濃度處處是 0。$r_0=1$、$R=4$：鏡像在 0.25、強度 $-0.25$。
由高斯定律，被球吃掉的比例是 $r_0/R=0.25$，逃到無窮遠的是 $(R-r_0)/R=0.75$：從距離 $R$ 出發的三維走路人，撞到半徑 $r_0$ 的球的機率是 $r_0/R$。

\subsubsection{電容的類比}
\bk{pp.~90--91}\ \anim{W407}

靜電學（高斯單位）：$\nabla^2\phi=-4\pi\rho$，導體上的電荷由高斯定律算；把它翻成擴散的語言（Fig.~4.7：內導體＋遠處的球殼）：
\[Q=-\frac1{4\pi}\int\nabla\phi\cdot d\vec S,\qquad \frac{-\int\vec j\cdot d\vec S}{c_\infty}=4\pi\beta C,\qquad C_{\rm sphere}=r_0.\tag{4.64--4.66}\]
讀法：\textbf{任何形狀的吸收體，每單位時間吃掉的量＝$4\pi\beta c_\infty\times$（它的電容）}。球的 $C=r_0$ 回到 (4.63)。
\begin{errbox}[p.~91，(4.64) 前]
書上寫 $\rho(\vec r)=-4\pi\nabla^2\phi(\vec r)$。高斯單位的 Poisson 方程是 $\nabla^2\phi=-4\pi\rho$，即 $\rho=-\frac1{4\pi}\nabla^2\phi$（(4.64) 本身是對的）。
\end{errbox}

\subsubsection{布滿受器的球}
\bk{pp.~91--93}\ \anim{W408}

\textbf{問}：細胞表面只有一小部分是受器（吸收點）。吃營養的速率會打幾折？

\textbf{書上的論證}：$n$ 個半徑 $a$ 的受器，每個帶電荷 $Q/n$。一個受器上的電位 $\approx\frac{Q}{na}$（自己）$+\frac{Q(n-1)}{nr_0}$（其他受器，像均勻帶電的球）。
$na\gg r_0$ 時第二項主導，電容 $\approx r_0$：和整顆球都是吸收面一樣好——即使受器只占表面的一小部分（$na^2\ll r_0^2$）。
\begin{errbox}[p.~92--93]
p.~92 的「$na\ll r_0$」（受器只占一小部分表面）：面積比是 $na^2/4r_0^2$，條件應為 $na^2\ll r_0^2$（即 $a\ll r_0/\sqrt n$），否則和 p.~93 的 $na\gg r_0$ 矛盾。
p.~93 的「$a\ll r/\sqrt n$」應為 $a\ll r_0/\sqrt n$。
\end{errbox}
\textbf{書外補充（Berg \& Purcell 1977, \emph{Biophys.\ J.} 20, 193）}：把受器看成嵌在不吸收表面上的圓盤（只有朝外的一面，等效電容 $a/\pi$；整個薄圓盤是 $2a/\pi=0.6366\,a$）。
$n$ 個受器並聯是 $na/\pi$，再和「從遠處走到細胞」的 $r_0$ 串聯：
\[\frac1C=\frac1{r_0}+\frac{\pi}{na},\qquad \frac{J}{J_{\max}}=\frac{C}{r_0}=\frac{na}{na+\pi r_0},\qquad n_{1/2}=\frac{\pi r_0}{a}.\]
例：細菌 $r_0=1\,\mu$m、受器 $a=1$ nm：$n_{1/2}=3142$ 個受器就達到一半的效率，它們只占表面的 0.079\%。
\begin{table}[H]
\centering
\small
\FourTabBP
\caption{Berg--Purcell（書外）：受器只占表面百分之幾，吸收就接近整顆球（$J_{\max}=4\pi\beta c_\infty r_0$）}\label{four-tab-bp}
\end{table}

% =====================================================================
\subsection{補充：吸收邊界}
\bk{pp.~93--94}\ \anim{W408}

格子上「牆到底在哪裡」要小心。取格點 $n\ge0$、懸崖緊鄰在 $0$ 的左邊（$-1$）。從 $x$ 走到 $y$ 的生成函數：
\begin{gather*}
G(z;x,y)=z\big(G(z;x,y-\Delta)+G(z;x,y+\Delta)\big)+\delta_{x,y}\quad(y\ge\Delta),\tag{4:S-1}\\
G(z;x,0)=zG(z;x,\Delta)+\delta_{x,0},\tag{4:S-2}\\
G(z;x,-\Delta)=0.\tag{4:S-3}
\end{gather*}
$y=0$ 少了左邊的鄰居，(4:S-2) 長得不一樣。技巧：加一個虛擬格點 $-1$，規定 (4:S-3)，(4:S-1) 就對所有 $y\ge0$ 成立。這就是\textbf{吸收邊界條件}：濃度為 0 的地方是「第一個不存在的格點」$-1$，不是 0。
\begin{errbox}[p.~93，(4:S-1)(4:S-2)]
(4:S-2) 印成 $G(z;0,y)=zG(z;x,\Delta)+\delta_{x,\Delta}$，應為 $G(z;x,0)=zG(z;x,\Delta)+\delta_{x,0}$。
(4:S-1) 後的「walks starting at $y$ and ending at $y$」應為「starting at $x$ and ending at $y$」。
\end{errbox}
\textbf{算到底}：牆在 $-1$，鏡像點是 $-x-2$（對 $-1$ 鏡射）：$C'(N;x,y)=C(N;y-x)-C(N;y+x+2)$。從 0 出發、$2k$ 步回到 0：
$\binom{2k}{k}-\binom{2k}{k+1}$＝Catalan 數 $1,\ 1,\ 2,\ 5,\ 14,\ 42$（$k=0,\dots,5$），和逐步遞迴一致。生成函數 $G'(z;x,y)=\frac{\lambda^{|x-y|}-\lambda^{x+y+2}}{\sqrt{1-4z^2}}$（$\lambda=\frac{1-\sqrt{1-4z^2}}{2z}$）；$z\to\frac12$ 時 $\lambda\to1$，得到從 $x$ 出發、掉下去之前到訪 $y$ 的平均次數 $2(\min(x,y)+1)$（例如 $x=1$、$y=3$ 是 4 次）：
正是 4.2.1 節的帳篷在 $L\to\infty$ 的極限。

% =====================================================================
\subsection{練習 4.1--4.6 解答}
{\small
\begin{enumerate}[label=\textbf{4.\arabic*},itemsep=0.05em]
\item 半空間 $x>0$（吸收牆在 $x=0$）的生成函數。\quad 答：鏡像，$G'(z;x_1,x_2)=G(z;x_1-x_2)-G(z;x_1+x_2)$（只動 $x$ 座標）。
一維格子：$G(z;d)=\lambda^{|d|}/\sqrt{1-4z^2}$，所以 $G'=\frac{\lambda^{|x_1-x_2|}-\lambda^{x_1+x_2}}{\sqrt{1-4z^2}}$，$\lambda=\frac{1-\sqrt{1-4z^2}}{2z}$（$z=0.3$ 時和逐步計數的級數一致）。
\item 由生成函數推 (4.2)(4.3)。\quad 答：三維連續版 $[z^N]\frac{1}{1-z(6-\Delta^2k^2)}=(6-\Delta^2k^2)^N\approx6^Ne^{-Nk^2/6}$，反傅立葉得 $6^N(\frac{2\pi N}{3})^{-3/2}e^{-3r^2/2N}$；
起點與鏡像點相減就是 (4.2)（$\mu=6$）。對 $y,z$ 積分，每個方向 $\int e^{-3y^2/2N}dy=(\frac{2\pi N}{3})^{1/2}$，兩個方向吃掉 $(\frac{2\pi N}{3})^{-1}$，剩 (4.3) 的 $(\frac{2\pi N}{3})^{-1/2}$。
\item 兩面牆的鏡像點。\quad 答：對 0 與對 $L$ 輪流鏡射：$x_0\to-x_0\to2L+x_0\to-2L-x_0\to\cdots$。正鏡像在 $2kL+x_0$、負鏡像在 $2kL-x_0$（$k\in\mathbb Z$）；
$C_{\rm box}=\sum_k[C(N;x-x_0-2kL)-C(N;x+x_0-2kL)]$。$L=6$、$x_0=3$、$N=10$、$x=3$：$252-45-45=162$。
\item 驗證正弦模態、推 (4.12)。\quad 答：$\sin q(x+\Delta)+\sin q(x-\Delta)=2\cos q\Delta\,\sin qx$，所以 $\sin qx$ 是「鄰居相加」的特徵向量，$q=n\pi/L$ 讓它在兩面牆上都是 0。
再用正交性 $\sum_{x=1}^{L-1}\frac2L\sin\frac{m\pi x}{L}\sin\frac{m'\pi x}{L}=\delta_{mm'}$ 把 $\delta_{x_1,x_2}=\sum_m\frac2L\sin qx_1\sin qx_2$ 展開，代進 (4.11) 比較 $\sin qx_1$ 的係數，就是 (4.12)。
\item 立方體裡的漫步。\quad 答：見 4.1.4 節與表 \ref{four-tab-shapes}：短時間 $6^NL^3$，長時間 $6^N\frac{512L^3}{\pi^6}e^{-\pi^2\Delta^2N/2L^2}$（$512/\pi^6=0.533$）；
格子 $L=6$：$N=0$ 是 125、$N=2$ 是 2970。
\item 被左牆吸收的機率。\quad 答：$P(x_0)=\frac{L-x_0}{L}$。由 (4.49)，流進左牆 $\beta c'(0)=s_0\frac{L-x_0}{L}$、右牆 $s_0\frac{x_0}{L}$；
也可以直接解 $P(x)=\frac12[P(x-1)+P(x+1)]$、$P(0)=1$、$P(L)=0$（直線）。$L=5$、$x_0=2$：$\frac35$（右牆 $\frac25$）；$L\to\infty$ 時 $\to1$。
\end{enumerate}}

% =====================================================================
\subsection{第 4 章的回顧與筆誤一覽}
\textbf{回顧}：吸收牆就是接地導體——鏡像、模態、電容，靜電學的工具全都能拿來數走法。
\begin{center}
\small\setlength{\tabcolsep}{4pt}
\begin{tabular}{@{}llll@{}}
\toprule
情況 & 工具與結果 & 書上的式子 & 動畫\\
\midrule
吸收牆 & 鏡像：碰過牆的路線相減 & (4.1)、(4.57)、(4:S-3) & W401、W406、W408\\
有限區域 & 正弦／$j_0$ 模態；長時間 $e^{-\pi^2N/2L^2}$ & (4.15)、(4.24)、(4.37) & W402--W404\\
穩定態 & Poisson：帳篷、$1/r$、偶極 & (4.41)、(4.49)、(4.52)、(4.58) & W405--W406\\
吸收體＝導體 & 被吃掉的量＝$4\pi\beta c_\infty\times$ 電容 & (4.63)、(4.66) & W407--W408\\
\bottomrule
\end{tabular}
\end{center}
\textbf{筆誤一覽}（給 2004 年版）：每一條都用程式核對過（\texttt{walk\_bc.check()}、\texttt{check\_numbers\_ch4()}），逐式的 LaTeX 版見 \texttt{ref/md/ch04\_formulas.md}。
(4.15) 的和與 p.~72 的 $x_0^2\gg2N$ 是書上沒處理好的地方（〔注〕），其餘在文中用紅框。
\begin{center}
\footnotesize
\FourTabErrata
\end{center}
